% Source: Wald GR, physical PDF p. 208
%         (printed p. 195).
% Coverage: A sequence in the chronological future of a closed set converging to its boundary.
% Figures/tables/footnotes: Figure 8.4; no tables or footnotes.
% Status: source-reviewed and compile-clean.
% Uncertainties: none.

\begingroup
\renewcommand{\thefigure}{8.4}
\begin{figure}[H]
  \centering
  \begin{tikzpicture}[x=.82cm,y=.82cm,>=Latex,line cap=round,font=\small]
    \draw[thick] (-2.55,.10) .. controls (-1.55,.58) and (-.42,.10) .. (.90,.20)
      -- (2.72,2.42);
    \node at (1.16,1.60) {\(I^+(C)\)};
    \draw[->] (.22,-.08) -- (.14,.32) node[below=.15cm] {\(C\)};
    \fill (-1.98,2.30) circle (1.6pt) node[left] {\(p\)};
    \foreach \x/\y in {-1.72/2.64,-1.57/2.98,-1.42/3.33,-1.30/3.65} {
      \fill (\x,\y) circle (1.25pt);
    }
    \node at (-1.55,3.92) {\(q_n\)};
    \draw[thick] (-2.55,.10) .. controls (-2.65,1.32) and (-2.33,2.26) .. (-1.72,2.64);
    \draw[thick] (-2.39,.14) .. controls (-2.34,1.50) and (-2.08,2.42) .. (-1.57,2.98);
    \draw[thick] (-2.22,.16) .. controls (-2.08,1.65) and (-1.80,2.70) .. (-1.42,3.33);
    \draw[thick] (-2.03,.18) .. controls (-1.82,1.75) and (-1.59,2.93) .. (-1.30,3.65);
    \draw[->] (-1.10,2.42) -- (-1.78,2.12);
    \node at (-.67,2.54) {\(\lambda_n\)};
  \end{tikzpicture}
  \caption{A spacetime diagram showing a sequence of points in \(I^+(C)\)
  converging to \(p\in\partial I^+(C)\) (see theorem~\ref{thm:8.1.6}).}
  \label{fig:8.4}
\end{figure}
\endgroup
